% TOP_PROBLEM: {"id": 3367, "title": "Schanuel's Conjecture", "category_id": 1, "category": {"id": 1, "name": "number_theory", "display_name": "Number Theory", "description": "Properties of integers, prime numbers, Diophantine equations.", "slug": "number-theory", "order_index": 1, "created_at": "2026-07-31T15:26:25.670Z"}, "status": "open"}
% FIRST_POSTED: 2026-09-25
% ENTRY_KIND: statement_only
% !TeX program = lualatex
% Definition and sources only; this file is not an LLM solution attempt.
\documentclass[11pt]{article}
\usepackage[margin=25mm]{geometry}
\usepackage{fontspec}
\setmainfont{Latin Modern Roman}
\usepackage{amsmath,amssymb,mathtools}
\usepackage{unicode-math}
\setmathfont{Latin Modern Math}
\usepackage{xurl,hyperref}
\hypersetup{colorlinks=true,urlcolor=blue}
\setcounter{secnumdepth}{0}
\setlength{\parindent}{0pt}
\setlength{\parskip}{5pt}
\setlength{\emergencystretch}{3em}
\begin{document}
\section*{\texorpdfstring{Schanuel's Conjecture}{Schanuel's Conjecture}}\label{3367}
\subsection{Definitions and mathematical statement}
Numbers $z_1,\ldots,z_n$ are ${\mathbb{Q}}$-linearly independent when $\sum_jq_jz_j=0$ with $q_j\in{\mathbb{Q}}$ forces all $q_j=0$. The transcendence degree of a field extension is the maximum size of an algebraically independent subset. Schanuel's assertion is
\[
 \dim_{{\mathbb{Q}}}\operatorname{span}_{{\mathbb{Q}}}\{z_1,\ldots,z_n\}=n
 \quad\Longrightarrow\quad
 \operatorname{trdeg}_{{\mathbb{Q}}}{\mathbb{Q}}(z_1,\ldots,z_n,e^{z_1},\ldots,e^{z_n})\ge n.
\]
It is quantified over every $n\ge1$ and every such complex tuple; the $z_j$ themselves need not be algebraic.
\par\smallskip\textit{Formulation and concept references:} \href{https://www.mdpi.com/2227-7390/9/7/717}{[S1]}.

\subsection{Short English statement}
Do linearly independent complex numbers, together with their exponentials, always contain at least as much algebraic independence as Schanuel predicts?
\subsection{Sources}
\begin{itemize}
\item[S1]Schanuel's Conjecture and the Transcendence of Power Towers.
\url{https://www.mdpi.com/2227-7390/9/7/717}.
\textit{Formulation source.}
\item[S2]Combining the conjectures of Schanuel and Zilber-Pink.
\url{https://content.ems.press/assets/public/full-texts/serials/rlm/36/3/14299588/online/10.4171-rlm-1086.pdf}.
\textit{Further reference; not a proof endorsement.}
\item[S3]Schanuel Property for Elliptic and Quasi-Elliptic Functions.
\url{https://webusers.imj-prg.fr/~michel.waldschmidt/articles/pdf/SchanuelPropertyAlmostAll.pdf}.
\textit{Further reference; not a proof endorsement.}
\end{itemize}
\end{document}
